% BHU PYQ -- Manually transcribed & error-corrected
% Paper: BCH-122 : Financial Accounting - II
% Exam: B.Com. (Hons.) Semester II Examination, 2022-23

\documentclass[11pt,a4paper]{article}
\usepackage[a4paper,top=2.5cm,bottom=2.5cm,left=2.8cm,right=2.8cm,headheight=14pt]{geometry}
\usepackage{amsmath,amssymb}
\usepackage[T1]{fontenc}
\usepackage[utf8]{inputenc}
\usepackage{lmodern,microtype,fancyhdr,enumitem,hyperref,lastpage,parskip}

\hypersetup{
    colorlinks=true,
    urlcolor=black,
    linkcolor=black,
    pdftitle={BCH-122: Financial Accounting - II -- BHU B.Com. Sem II 2022-23}
}

\pagestyle{fancy}
\fancyhf{}
\renewcommand{\headrulewidth}{0.4pt}
\renewcommand{\footrulewidth}{0.4pt}
\fancyhead[L]{\small\textit{Banaras Hindu University}}
\fancyhead[R]{\small\textit{BCH-122: Financial Accounting - II}}
\fancyfoot[C]{\small Page \thepage\ of \pageref{LastPage}}

\newlist{parts}{enumerate}{2}
\setlist[parts,1]{label=(\alph*),leftmargin=*,itemsep=5pt,topsep=3pt}
\setlist[parts,2]{label=(\roman*),leftmargin=*,itemsep=3pt,topsep=2pt}
\newcommand{\pts}[1]{\hfill\textit{[#1]}}

\begin{document}

\begin{center}
  {\large\bfseries BANARAS HINDU UNIVERSITY}\\[2pt]
  {\normalsize B.Com. (Hons.) --- Semester II Examination, 2022-23}\\[6pt]
  \rule{0.8\linewidth}{0.5pt}\\[6pt]
  {\large\bfseries Paper No. BCH-122: Financial Accounting - II}\\[4pt]
  {\normalsize Time: 3 Hours\hfill Maximum Marks: 70}
\end{center}

\rule{\linewidth}{0.4pt}

\smallskip
\noindent\textit{Instructions: Answer all questions. Marks are indicated against each question.}

\bigskip

\noindent\textbf{Question 1.}
From the following information prepare a Bank Reconciliation Statement for the year ended 30th September 2022:\pts{14}
\begin{center}
  \small
  \begin{tabular}{|l|r|}
    \hline
    \textbf{Particulars} & \textbf{Amount (Rs.)} \\
    \hline
    Balance as per Cash Book & 12,000 \\
    Cheques issued but not encashed & 8,000 \\
    Cheques deposited but not credited by bank & 7,600 \\
    Bank credited interest & 1,200 \\
    Bank debited Insurance Premium & 1,500 \\
    A customer directly deposited Cash into Bank & 2,500 \\
    Cheque received entered into Cash Book but not sent to Bank & 900 \\
    Bank charges debited & 100 \\
    \hline
  \end{tabular}
\end{center}

\centerline{\textbf{OR}}

Differentiate between Bills of Exchange and Promissory Note. Show the journal entries in the books of the Drawer and the Drawee.\pts{14}

\medskip\rule{\linewidth}{0.2pt}

\noindent\textbf{Question 2.}
$X$ Company Ltd. issued a prospectus inviting applications for 2,000 shares of Rs. 10 each at a premium of Rs. 2 per share payable as:
\begin{itemize}[label=-,leftmargin=1.5cm]
  \item On Application: Rs. 2
  \item On Allotment: Rs. 5 (including premium)
  \item On First Call: Rs. 3
  \item On Second \& Final Call: Rs. 2
\end{itemize}
Applications were received for 3,000 shares and pro-rata allotment was made on applications for 2,400 shares. The remaining applications were rejected. Money overpaid on application was adjusted on allotment.

$R$, to whom 40 shares were allotted, failed to pay allotment money and on his subsequent failure to pay first call, his shares were forfeited. Mr. Mohan, the holder of 60 shares, failed to pay the two calls and his shares were forfeited.

All these forfeited shares were reissued as fully paid for Rs. 8 per share. Give the necessary journal entries.\pts{14}

\centerline{\textbf{OR}}

\begin{parts}
  \item Explain the conditions for the redemption of Preference Shares.\pts{6}
  \item Give the necessary journal entries for the redemption of Preference Shares.\pts{8}
\end{parts}

\medskip\rule{\linewidth}{0.2pt}

\newpage

\noindent\textbf{Question 3.}
Distinguish between a share and a debenture. Mention the types of shares.\pts{14}

\centerline{\textbf{OR}}

On 1st January 2018, $Z$ Ltd. issued 6\% debentures of Rs. 100 each to be redeemed at the end of four years. Sinking fund policy was followed to redeem the debentures. At the end of the 4th year, the investments were sold for Rs. 3,50,000 and the debentures were paid off. Sinking fund table shows that Rs. 0.228592 @ 6\% p.a. for 4 years provides Rs. 1. Give the necessary journal entries.\pts{14}

\medskip\rule{\linewidth}{0.2pt}

\noindent\textbf{Question 4.}
$B.W.$ Ltd. was incorporated on 1st March 2015 and received its certificate to commence business on 1st April 2015. It took over the business of $DK$ Brothers running from 1.1.2015. From the following particulars, calculate the profit prior to incorporation and post-incorporation in the books of $B.W.$ Ltd.:\pts{14}
\begin{parts}
  \item The Total sales for the year was Rs. 9,00,000, of which sales upto 31.03.2015 was only Rs. 3,00,000.
  \item The Trading Account disclosed a Gross Profit of Rs. 2,49,000.
  \item The Profit and Loss Account showed the following expenses:
  \begin{center}
    \small
    \begin{tabular}{|l|r||l|r|}
      \hline
      \textbf{Particulars} & \textbf{Amount (Rs.)} & \textbf{Particulars} & \textbf{Amount (Rs.)} \\
      \hline
      Salaries & 20,000 & Commission on Sales & 4,800 \\
      Discount on Sales & 3,600 & Auditor's fees & 12,500 \\
      Rent \& taxes & 28,000 & Establishment expenses & 11,400 \\
      Interest to Vendors & 10,000 & Depreciation & 12,000 \\
      Director's fees & 3,500 & Bad Debts (Rs. 500 pre-31st March) & 900 \\
      Printing \& Stationery & 5,200 & Advertising & 24,000 \\
      Interest on Debentures & 8,000 & Preliminary Expenses & 6,000 \\
      \hline
    \end{tabular}
  \end{center}
\end{parts}

\centerline{\textbf{OR}}

What is the basis of apportionment of pre and post incorporation profit? Explain and illustrate.\pts{14}

\medskip\rule{\linewidth}{0.2pt}

\newpage

\noindent\textbf{Question 5.}
Jagaran Ltd. had a capital of Rs. 6,00,000 divided into shares of Rs. 10 each. The following balances were available on 31st December 2022:\pts{14}
\begin{center}
  \small
  \begin{tabular}{|l|r||l|r|}
    \hline
    \textbf{Particulars} & \textbf{Amount (Rs.)} & \textbf{Particulars} & \textbf{Amount (Rs.)} \\
    \hline
    Opening Stock & 75,000 & Sales & 4,15,000 \\
    Purchases & 1,85,000 & Provision for Bad Debts & 3,500 \\
    Wages & 84,800 & Share Capital & 4,60,000 \\
    Freight & 13,100 & P \& L Balance (credit) & 14,500 \\
    General Expenses & 16,900 & Bills Payable & 38,000 \\
    Salaries & 14,500 & Sundry Creditors & 50,000 \\
    Interest on Debentures & 9,000 & 6\% Debentures & 3,00,000 \\
    Director's fees & 5,740 & Fixtures & 7,200 \\
    Preliminary Expenses & 5,000 & Cash at Bank & 39,900 \\
    6\% Debentures & 3,00,000 & Cash in hand & 750 \\
    Bad Debts & 2,110 & Premises & 3,00,000 \\
    Interim Dividend & 7,500 & Calls in Arrear & 7,500 \\
    Plant \& Machinery & 3,60,000 & 4\% Govt. Securities & 60,000 \\
    Goodwill & 25,000 & Sundry Debtors & 87,000 \\
    \hline
  \end{tabular}
\end{center}
Prepare Profit and Loss Account and Balance Sheet relating to 2022 from the above information after considering the following adjustments:
\begin{parts}
  \item Depreciate Plant \& Machinery by 10\% p.a. and fixture by 5\% p.a.
  \item Write off 1/5th of Preliminary expenses.
  \item Rs. 10,000 of wages were spent on erection of premises.
  \item Provide 5\% for Bad Debts on Debtors.
  \item Provide a final dividend @ 5\%.
  \item Transfer Rs. 10,000 to General Reserve.
  \item Provision for Income Tax Rs. 25,000.
  \item Closing Stock was valued at Rs. 1,01,000.
\end{parts}

\centerline{\textbf{OR}}

Give the proforma of Balance Sheet as per Indian Companies Act and explain it in detail.\pts{14}

\vfill
\rule{\linewidth}{0.4pt}
\begin{center}\small
\href{https://bhu-exam.vercel.app/}{Click here to visit our website}\\[4pt]
\href{https://me-aryan.vercel.app/}{Click here to visit our contributor}\\[4pt]
\href{https://quashan-me.vercel.app/}{Click here to visit our contributor}
\end{center}

\end{document}
